\section{Conclusion}
\label{sec:conclusion}

We have rigorously established the Squiggly Line Conjecture through formal mathematical proofs. Our fundamental theorem proves that complex curves systematically achieve lower expected distances to random points, with the key principle: greater complexity → greater arc length → greater covered volume → lower expected distance. We further demonstrated that distance reduction follows a submodular pattern, explaining the natural saturation of overparameterization benefits.

Our theoretical framework reveals dimension-independent scaling $\mathbb{E}[L(f_w, \mathcal{D})] = L_0 \cdot w^{-\alpha} + L_\infty$ across dimensions 2-10, with universal saturation at $w^* = \Theta(\sqrt{n})$. This geometric perspective provides unified explanations for the lottery ticket hypothesis, overparameterization benefits, and optimal architectural design, showing that early complexity additions provide disproportionately large benefits while later additions exhibit diminishing returns.

\textbf{Limitations and future work}: While effect sizes are typically 10-50\% improvements and benefits saturate around $w = \sqrt{n}$, future directions include architectural optimization for coverage, quantum machine learning applications, and biological validation.

\section*{Acknowledgments}

We thank the machine learning community for valuable discussions and insights that shaped this work.

\section*{Broader Impact}

This theoretical work may advance understanding of neural network initialization and overparameterization through geometric principles, with potential applications to more efficient architectures and training procedures.
